\section{Discussion}
\label{sec:discussion}

The operational detector's penultimate features strongly distinguish the
selected observed and injected populations, showing that its trained
representation can serve as a detector-specific diagnostic. The same
separation appears in self-supervised survey features, generic ImageNet
features, and elementary pixel statistics, so it is not confined to the
lens-trained representation.

The controls add two distinct observations: a leading pixel-statistic offset
predates insertion, and matched injections have a near-zero mean signed
movement along the separately fitted operational observed--bare axis, with
heterogeneous pairwise shifts. The separate recovery, source-photon,
source-construction, and fixed-target brightness experiments characterize
responses to controlled image changes. The matched controls locate differences
that predate injection, while the recovery experiments measure how the fixed
detector responds when source photons, source construction, or brightness are
changed.

\subsection{What the recovery measurements represent}
\label{sec:discussion_recovery}

The detector responds to 39/43 selected confirmed systems but to 3.45 per cent
of the fixed injection grid. These summaries describe different cohorts and
weighting schemes. The selected systems come from heterogeneous discovery
programmes, whereas \(G_{\rm rec}\) averages equally over nine observing
cells, 4,500 targets, and 11 Einstein radii
\citep{Huang2020,Huang2021,Stein2022,Storfer2024,Sonnenfeld2023,Sonnenfeld2024}.

On the 43 fixed targets, expectation-image recovery peaks at magnitude 20.5
and is lower at the two brighter nodes. Thus recovery is non-monotonic over
the tested intrinsic-magnitude range.

\subsection{Feature diagnostics, matching, and mass-model scope}
\label{sec:discussion_deflector}

The fixed lens head is one affine projection of the operational feature
vector. Its AUC of 0.980 shows that the existing score already expresses a
strong separating direction. Using the same cross-fitted readout for all four
representations shows whether the operational features are uniquely sensitive
to population identity. The probe and fixed head are complementary readouts
of the same feature vector, while the pixel-statistic AUC of 0.981 shows that
elementary image summaries also carry a strong separating direction.

The leading native-pixel offset is already present between the observed
references and bare targets (Section~\ref{sec:matched_probe_results}).
Foreground light, background estimation, survey processing, and sample
selection are all plausible contributors to that offset. The measured offset
is present before insertion; these analyses do not separate the possible
contributions listed above.

\(G_{\rm match}\) asks how one fixed injection moves each matched target along
an observed--bare axis. Its operational mean is near zero, but individual
shifts range from \(-0.453\) to \(+0.998\).
The near-zero cohort mean therefore coexists with substantial individual
movements. Matching reduces differences in the five recorded central-object
and observing quantities. The observed--bare readout describes differences
between complete cutouts; insertion measures movement along that readout on
the same targets. Redshift, velocity dispersion, stellar mass, source
properties, and environment remain outside that matching set. The radius
curve similarly describes response over assigned grid values rather than
target-specific inferred lens models.

\subsection{Arc-photon-noise treatment}
\label{sec:discussion_noise}

The negative aggregate photon effect persists across the 12
validation-selected checkpoints
(Section~\ref{sec:checkpoint_sensitivity_results}). The primary checkpoint's
radius-specific band is a different summary: it includes zero through
1.5 arcsec and lies below zero at larger tested radii
(Section~\ref{sec:poisson_results}).

Only the extra shot-noise term associated with injected flux is changed.
The observed target already contains its realised sky, detector, coaddition,
neighbours, and foreground light. The adopted effective gain is approximate,
and the intervention operates on post-coadd cutouts rather than reconstructing
exposure-level correlated noise or resampling. Eight draws control
finite-photon Monte Carlo error for the adopted
\(g=150\ {\rm e^{-}\,nmgy^{-1}}\) source-photon law. Gain-model and
threshold-calibration uncertainty remain outside the reported interval.

\subsection{Source-construction result}
\label{sec:discussion_source}

Under the tested construction, single-S\'{e}rsic recovery is about
0.56 percentage points higher than structured-source recovery. All 12
validation-selected checkpoint estimates have the same sign with
7/12 intervals excluding zero. Under the terminal-epoch policy all 12 estimates
retain that sign, while 8/12 intervals exclude zero. The direction is
therefore consistent across the examined policies, but interval exclusion is
checkpoint-policy-dependent. The estimand compares the two complete
source-construction pipelines.

The empirical arm is constructed from selected high-S/N F814W images by
deconvolution and reconvolution to a common resolution, followed by
non-negative projection, centring, support restriction, and normalization.
Its spatial map therefore contains both source structure and the effects of
measurement and regularization. Each arm reuses its own F814W-derived map
across \(g,r,z\); the tested construction has no colour gradients or
chromatic morphology.
Both arms share the selected targets, assigned SIE-plus-shear mappings,
effective PSF, modelled photon law, threshold, and post-coadd insertion.
In the noiseless known-source test, global flux, centring, and normalization
errors remain within the stated tolerances
(Appendix~\ref{app:source_transfer_followup}).

Future paired studies can vary source support, regularization, correlated
measurement noise, and
band-dependent morphology while continuing to hold targets, lens mappings,
fluxes, and photon draws fixed.

\subsection{Renderer qualification and numerical scope}
\label{sec:discussion_renderer}

Agreement in the fixed 142-scene validation cohort supports the numerical
implementation within its tested configurations
(Appendix~\ref{app:renderer_reconciliation}). These numerical checks do not
test the astrophysical adequacy of the adopted source, lens, PSF, or detector
models.

\subsection{Relation to published pipelines and resolution}
\label{sec:discussion_literature}

Published pipelines differ in their source constructions, deflector
information, and observing regimes: fully simulated \textit{Euclid}-like
images expose selection dependence on source and lens parameters
\citep{Herle2024}, whereas HSC and other recent pipelines use real galaxy
stamps and, in some cases, more tightly constrained deflector information
\citep{Canameras2021,Rojas2022,Schuldt2023,Canameras2024}. These approaches
vary choices beyond the paired structured--S\'{e}rsic intervention studied
here.

Resolution is especially important. At DR10 resolution, the effective PSF
smooths the rendered lensed images and reduces the observable fine structure
inherited from the source templates. HST, \textit{Euclid}, and seeing-limited
ground-based data resolve different source structure, so the magnitude and
direction of source-construction effects can vary with resolution. Source,
mass, PSF, detector, and processing choices should therefore be tested at the
resolution at which a selection function will be applied.

\subsection{Limitations and future work}
\label{sec:limitations}

Inference is restricted to the sampled systems and tested pipelines.
Generalizing these results to another lens population is limited by missing
physical covariates, mass parameters assigned independently of the target
galaxies, single-band source morphology, post-coadd insertion, approximate PSF
and gain models, and unknown overlap with Stein's pretraining data. The
operational evaluations retain the detector's trained 20--32-pixel annulus and
clip-10 map, making that mapping part of the operational estimand.

The detector panel contains 12 correlated variants of one architecture and
training pipeline. Sampler and schedule changes are
evaluated only at seed 42, the positive-example-set contrast changes several
properties together, and no between-checkpoint contrast interval is attached.
Each estimate conditions on an empirically calibrated threshold, so the panel
combines training variation with calibration-tail variation. It therefore
tests variation within this training family rather than across independent
detector architectures.
The operational feature analysis uses only the primary checkpoint. All four
probes share the 43 observed references; dependency-closed folds keep linked
images together, and component-equal weights prevent repeated injections from
receiving disproportionate weight.

Future work should first construct a deflector sample with reliable redshifts,
velocity dispersions or other mass proxies, then assign mass and environment
jointly with observed light. Controlled mass-model factors should include
density slope, composite baryonic and dark components, external convergence,
multiple deflectors, and line-of-sight structure. Each factor should be varied
within matched targets rather than inferred from an unpaired population
contrast.

Source tests should use multiple real-galaxy catalogues and latent-source
methods, band-dependent morphology and colour gradients, with support,
regularization, and correlated noise varied explicitly. The
structured--S\'{e}rsic intervention should be repeated at space-based and
ground-based resolution under the same underlying systems.
Exposure-level injection before coaddition would test PSF, noise correlation,
deblending, masking, and processing provenance that cutout-level addition
leaves fixed.

Future studies should predefine the simulated population, detector, and
thresholds; preserve system- or target-level grouping; report support in every
conditional cell; and combine unpaired comparisons with bare-target and paired
controls. Repeating the diagnostic under such controlled pipeline changes can
identify which changes alter observed--injected separation. Pairing those
changes with selection-function and downstream population-inference metrics
would connect the diagnostic to survey performance.
